% =============================================================================
\section{Easy\_SLM System Design}
% =============================================================================

Easy\_SLM is a browser-based single-page application (Next.js, React,
TypeScript, Zustand) that routes the same user specification to either a local
SLM via the effGen API~\cite{effgen} or a cloud LLM (OpenAI / Anthropic) and
returns streamed responses through a unified Server-Sent Event contract.
The chat layout---colors, composer position, message bubbles---is identical
across conditions; the only difference is the backend.
This parity is deliberate: it prevents interface aesthetics from confounding
model capability comparisons (RQ3, RQ4).

We derived four design goals from the formative task-selection study
(Section~\ref{sec:tasksel}) and the prior literature on Gulf of
Envisioning~\cite{subramonyam2024gulf, zamfirescu2023johnny}.

% -----------------------------------------------------------------------------
\subsection{DG1: Paired SLM--LLM Execution with Objective Scoring}
% -----------------------------------------------------------------------------

\textit{Rationale.}
RQ1 and RQ2 require (i)~matched inputs submitted to both backends under
the same specification and (ii)~automatic metric computation.
Letting users freestyle prompts in separate tools confounds model capability
with prompting skill; a single shared interface eliminates that confound.

\textit{Implementation.}
The study home page presents three task cards, each with two entry buttons---
\textit{Begin with SLM} and \textit{Begin with LLM}---so participants explicitly
choose a condition per task.
A mode badge (CPU icon / Cloud icon) persists in the task bar throughout the
session.
Completion is tracked independently per task and mode: submitting the SLM
attempt does not mark the LLM attempt done, and vice versa.
The backend logs provider, model, task, and full prompt for post-hoc metric
computation (ROUGE, F1, rule-based pass rates) against human-written
references without exposing scores during the task.

% -----------------------------------------------------------------------------
\subsection{DG2: Translate User Intentions into Concrete Goals}
% -----------------------------------------------------------------------------

\textit{Rationale.}
Novice users tend to underspecify intent, and SLMs are especially brittle to
vague input~\cite{zamfirescu2023johnny}.
The \textit{intentionality gap}---not knowing what to ask---and the
\textit{capability gap}---not knowing what the model can produce---both
contribute to the Gulf of Envisioning~\cite{subramonyam2024gulf}.

\textit{Implementation.}
A two-step \textbf{Goal Wizard} modal guides users through (1)~writing a
one-sentence goal with an optional reframing prompt (e.g., ``turn my idea into
a clear, step-by-step task''), and (2)~selecting audience (\textit{just for me
/ team / general public}), output format (\textit{paragraph / checklist / table
/ JSON}), constraints, and success criteria.
In the SLM condition the wizard pre-loads a task-appropriate example goal that
participants can edit; in the LLM condition participants go directly to chat, so
the effect of goal scaffolding is isolatable.
The committed intent is compiled into a system prompt and displayed in a
persistent \textbf{Intent Ribbon} confirming the specification is live before
the first message is sent.

% -----------------------------------------------------------------------------
\subsection{DG3: Scaffold Prompt Design}
% -----------------------------------------------------------------------------

\textit{Rationale.}
Even after goal clarification, translating intent into effective instructions
remains hard---the \textit{instruction gap}.
Natural-language ambiguity and ``prompt brittleness'' cause SLM outputs to
vary unpredictably~\cite{zamfirescu2023johnny}, a confound that structured
decomposition can reduce.

\textit{Implementation.}
The \textbf{Prompt Plan Editor} lets users build an ordered list of prompt
steps.
\textit{Decompose with effGen} sends the current prompt text to effGen's
rule-based decomposition engine, returning subtasks that replace the step list
(no external LLM required for decomposition).
A \textbf{Variant Generator} produces three prompt variants (strict /
brief / explicit) side-by-side so users can inspect how phrasing changes
output before committing.
A \textbf{Template Gallery} offers four pre-built scaffolds---checklist,
summarize for audience, extract to JSON, rewrite with constraints---lowering
the entry floor for first-time users.
The compiled prompt is always visible in an accordion preview so users can
confirm what will be sent.

% -----------------------------------------------------------------------------
\subsection{DG4: Guided Model Selection}
% -----------------------------------------------------------------------------

\textit{Rationale.}
The SLM ecosystem offers hundreds of models varying in quantization, context
length, and hardware requirement.
Novice users face choice overload and lack heuristics for matching a model
to their task and device~\cite{subramonyam2024gulf}.

\textit{Implementation.}
The \textbf{Right Inspector} panel recommends models from a curated catalog
based on the committed intent and estimated compute tier.
A \textit{Reuse loaded model} toggle prevents costly model reloads across
turns.
In the LLM condition the same panel slot shows a provider selector
(OpenAI / Anthropic) and model dropdown, preserving layout parity while
removing SLM-specific guidance---so any right-panel effects measured in RQ4
are attributable to content, not position.
